% =====================================================================
%  Chapter 16 — Non-implication (S44)
% =====================================================================
\chapter{Two Hypotheses That Are Not One}
\label{ch:non-implication}

\begin{record}
\textbf{What this chapter covers.} Session \Sn{44}. The proof document had
three times described \texttt{conj:K-refined} (the \(c_K\) decorrelation
hypothesis, \Sn{34}) and \texttt{target:band-absorption} (the annulus
replacement target, \Sn{41}) as being ``of the same analytic type'', and had
observed that their crude bounds fail by exactly the same single logarithm.
Nobody had checked whether the resemblance was substantive. If one implied the
other, or both followed from a common hypothesis, the programme's open-item
count would fall --- the same kind of win as \Sn{37}'s unification.

\smallskip
The answer is no, in a precisely bounded sense, and this chapter is careful
about the bound because the programme's own summaries were not. Two things in
particular must be stated exactly, because both were stated wrongly once and
corrected: the direction of the inequality relating the two statistics, and
the scope of the word ``independent''.
\end{record}

\section{The setting: what a non-implication can even mean}
\label{sec:s44-setting}

Neither conjecture can be refuted or verified by exhibiting a Navier--Stokes
solution. Both are asserted only \emph{near a putative Type-I singularity},
whose existence is the open question. So the question \Sn{44} could actually
settle is narrower and is stated in the document before any result:
\emph{does either implication follow from the a priori structure the two
hypotheses are formulated against?}

\begin{defn}[Admissible test configuration --- \texttt{def:s44-config}, line~8835]
\label{def:s44config}
Fix \(\nu>0\), \(M\ge1\), \(\rstar=M^{-1/2}\), \(\Lfac\ge1\) and a cutoff
\(\psi\). An admissible test configuration is a triple \((m,w,h)\) on
\(Q(\rstar)\) with \(m\) Lipschitz and \(w,h\ge0\) measurable, satisfying
\begin{enumerate}[label=\textup{(A\arabic*)},leftmargin=2.6em]
\item \(0\le m\le M\) and \(\sup_{Q(\rstar)}m=M\);
\item the enstrophy budget
  \(\nu\iint_{Q(\rstar)}\left(\abs{\nabla m}^2+m^2w\right)\le C_{\mathcal E}\Lfac M^{1/2}\);
\item the pointwise link \(\abs{\nabla w}\le2w^{1/2}h^{1/2}\), with \(h\)
  standing for \(\abs{\nabla^2\ehat}^2\).
\end{enumerate}
\end{defn}

That is the whole of the structure both hypotheses are phrased in. Nothing
else is imposed: no evolution equation, and no requirement that \(w\) actually
be \(\abs{\nabla\ehat}^2\) for a unit field.

\section{The structural reason, found first}
\label{sec:disjoint}

The result that makes everything else obvious once stated is a remark about
supports.

\begin{statement}[\texttt{rem:s44-disjoint}, line~8771]
\label{st:disjoint}
\texttt{conj:K-refined} is a statement about \(E=Q(\rstar)\cap\{m\ge M/4\}\).
\texttt{target:band-absorption} is a statement about
\(\Band\subseteq\{m\le M/4\}\). The two regions meet only in the level set
\(\{m=M/4\}\), which carries no measure whenever \(M(t)/4\) is a regular
value. \textbf{Neither hypothesis constrains any integral supported in the
other's region.}
\end{statement}

``Same analytic type'' is a statement about the \emph{form} of the two
hypotheses --- both are decorrelation statements normalised the same way ---
and not about their content. They are the same kind of statement about
disjoint parts of the same object.

\begin{figure}[ht]
\centering
\begin{tikzpicture}[font=\scriptsize]
% vertical axis of |omega| levels
\draw[thick,->] (0,-0.3) -- (0,5.4) node[above] {\(\abs{\omega}\)};
\foreach \y in {5.0, 2.9, 1.9, 1.25}{ \draw (-0.1,\y) -- (0.1,\y); }
\node[left] at (-0.15,5.0) {\(M\)};
\node[left] at (-0.15,2.9) {\(M/4\)};
\node[left] at (-0.15,1.9) {\(M/8\)};
\node[left] at (-0.15,1.25) {\(M/16\)};
% E region
\fill[provedgreen, opacity=0.18] (0.5,2.9) rectangle (5.3,5.0);
\draw[provedgreen, thick] (0.5,2.9) rectangle (5.3,5.0);
\node[provedgreen, font=\small\bfseries] at (2.9,4.5) {\(E\)};
\node[align=center, text width=44mm] at (2.9,3.55)
  {\(E=Q(\rstar)\cap\{m\ge M/4\}\)\\
   \texttt{conj:K-refined} lives here: \(c_K\le C/\Lfac\)};
% Band region
\fill[impossiblered, opacity=0.16] (0.5,1.9) rectangle (5.3,2.9);
\draw[impossiblered, thick] (0.5,1.9) rectangle (5.3,2.9);
\node[impossiblered, font=\small\bfseries] at (1.15,2.62) {\(\Band\)};
\node[align=center, text width=42mm] at (3.2,2.28)
  {\texttt{target:band-absorption} lives here};
% overlap line
\draw[very thick, black] (0.5,2.9) -- (5.3,2.9);
\node[right, align=left, text width=32mm] at (5.45,2.9)
  {they meet only on \(\{m=M/4\}\):\\ a null set};
% LC region bracket
\draw[decorate, decoration={brace, amplitude=5pt}, openblue, very thick]
  (5.9,1.25) -- (5.9,5.0);
\node[right, openblue, align=left, text width=34mm] at (6.25,3.1)
  {\textbf{(LC)}, \texttt{conj:lc-s44}:\\ verbatim the \(c_K\) hypothesis,
   asserted on the \emph{larger} region \(\{m\ge M/16\}\) --- which contains
   both};
\draw[openblue, dashed] (0.5,1.25) rectangle (5.3,5.0);
\end{tikzpicture}
\caption[Why neither hypothesis sees the other]{The level decomposition of
\texttt{rem:s44-disjoint} (line~8771). The two hypotheses are supported on
disjoint regions, so neither constrains the other, and the counterexamples of
\S\ref{sec:counterexamples} merely make that concrete. (LC) covers both by
being asserted on a region containing both --- which is why it implies neither
of them as a statement.}
\label{fig:s44-levels}
\end{figure}

The disjointness has independent corroboration from the programme's own code.
\texttt{layer4/test\_band\_concentration.py}, written at \Sn{42} with none of
this in view, contains a test named
\texttt{test\_band\_and\_E\_are\_essentially\_disjoint}, asserting exactly this
numerically.

\section{The two non-implications}
\label{sec:counterexamples}

\begin{statement}[\texttt{prop:s44-b-fails}, line~8879]
\label{st:bfails}
There is a family of admissible test configurations in which both band
left-hand sides vanish identically --- so
\texttt{target:band-absorption} holds with \(C=0\) --- while
\(c_K\ge\abs{E}/(8\ell^5)\to\infty\). Band absorption therefore implies
\emph{no upper bound on \(c_K\) whatsoever}.
\end{statement}

\begin{statement}[\texttt{prop:s44-a-fails}, line~8980]
\label{st:afails}
There is a family with \(c_K=0\) --- the strongest possible form of
\texttt{conj:K-refined} --- for which every admissible level pair of width
\(\ge1/16\) violates \eqref{eq:bandY} by a factor \((\rstar/\ell)^5\to\infty\).
The adversary is \(\abs{\omega}\) oscillating rapidly \emph{inside} the band,
so \(\abs{\nabla m}^2\) and \(w\) co-concentrate exactly where the \(c_K\)
hypothesis has no purchase.
\end{statement}

\section{The scope, stated as the document states it}
\label{sec:s44-scope}

This is the first of the two things that must be exact.

\begin{record}
\textbf{\texttt{rem:s44-scope}} (line~8855), quoted. ``An admissible test
configuration need not come from a solution of the Navier--Stokes equations:
\(w\) need not equal \(\abs{\nabla\ehat}^2\) for a unit field \(\ehat\), \(h\)
need not be its second-derivative density, and no evolution equation is
imposed. Consequently \texttt{prop:s44-b-fails} and \texttt{prop:s44-a-fails}
do \emph{not} say that the two conjectures are independent as statements about
Navier--Stokes solutions; that question is not decidable by any argument
available here. What they do say, and all they say, is: \emph{no implication
between} \texttt{conj:K-refined} \emph{and} \texttt{target:band-absorption}
\emph{can be derived from} (A1)--(A3), i.e.\ from the a priori budgets and the
pointwise structure in terms of which the two statements are phrased.''
\end{record}

\begin{obsv}[``Proved independent'' is the wrong phrase, and the programme used it]
\label{obs:not-independent}
\texttt{PROGRESS.md}'s \Sn{44} entry is headed ``\(c_K\) and band-absorption
proved INDEPENDENT''. The session log's own summary uses ``proved independent
rather than suspected equivalent''. Both drop the qualifier, and
\texttt{rem:s47-handover-corrigenda} in the walls document records the
correction: ``\emph{Proved independent} without that qualifier overstates
them.''

\smallskip
The distinction is not pedantry. An implication between the two conjectures
could still hold for Navier--Stokes solutions, supplied by PDE structure that
\Sn{44}'s test class deliberately omits; what is ruled out is deriving one
from the other using only the budgets both are written against. A future
session that reads ``independent'' and stops looking would be acting on a
statement the document does not support.
\end{obsv}

\section{The inverted lemma}
\label{sec:s44-inversion}

This is the second thing that must be exact, and it is the part of the session
most worth reading.

Round one of the delegated work returned a favourable headline:
\texttt{target:band-absorption} is strictly \emph{weaker} than
\texttt{conj:K-refined}, needing only an \(O(1)\) equipartition constant where
\(c_K\) needs the decaying \(C/\Lfac\); and therefore the programme's four
sessions of measurements at \(c_K\approx1\) already supported it directly.
That would have been genuinely good news, and it would have redirected later
sessions toward the wrong target.

The supporting lemma was inverted.

\begin{statement}[\texttt{lem:s44-xi-vs-cK}, line~8700 --- the correct direction]
\label{st:xi-vs-ck}
With \(\Xi\) the volume-normalised coupling excess and \(c_K\) as in
\texttt{conj:K-refined},
\[
  \Xi[\mathbf 1_E] \;=\; c_K\,\frac{\abs{Q(\rstar)}}{\abs{E}} \;\ge\; c_K ,
\]
and more generally \(\Xi[\varpi]=c_K[\varpi]\,\abs{Q(\rstar)}/\abs{R}\) for any
weight \(\varpi\) supported in \(R\subseteq Q(\rstar)\). Consequently:
\begin{enumerate}[label=(\roman*),leftmargin=2.4em]
\item a bound \(\Xi[\varpi]\le\kappa\) is \emph{strictly stronger} than the
  correlation bound \(c_K[\varpi]\le\kappa\), by exactly the volume ratio;
\item \texttt{conj:K-refined} gives only
  \(\Xi[\mathbf 1_E]\le(C/\Lfac)\abs{Q(\rstar)}/\abs{E}\), which is \(O(1)\)
  \emph{only} under a volume lower bound on \(E\);
\item \(\Xi[\varpi]\ge1\) whenever the two densities share a support of measure
  \(\le\abs{Q(\rstar)}\) and are uncorrelated on it --- constants on \(R\) give
  \(c_K=1\) but \(\Xi=\abs{Q(\rstar)}/\abs{R}\gg1\).
\end{enumerate}
\end{statement}

The first draft had recorded \(\Xi[\mathbf 1_E]=c_K\abs{E}/\abs{Q(\rstar)}\le c_K\)
--- the reciprocal volume factor, hence the wrong direction.

\begin{record}
\textbf{How it was caught}, from \texttt{rem:s44-corrigendum} (line~8741) and
the session log. Three checks, all cheap:
\begin{enumerate}[label=(\roman*),leftmargin=2.4em]
\item re-derive the algebra directly from \texttt{conj:K-refined}'s own
  definition and \texttt{def:s44-excess};
\item a numerical check on four random configurations --- the stated form
  matched nothing; the corrected form matched to machine precision;
\item the sanity case: \(\abs{\nabla m}^2\) and \(w\) constant on \(E\) and
  zero off it give \(c_K=1\), uncorrelated as they should be, but
  \(\Xi=\abs{Q}/\abs{E}\gg1\).
\end{enumerate}
The error also contradicted \emph{two of the same submission's own results}
--- \texttt{rem:s44-jle-status}, which correctly computes
\(\Xi=\abs{Q(\rstar)}/\abs{S}\) under concentration, and
\texttt{prop:s44-transport} --- and that internal inconsistency is what drew
attention to it.
\end{record}

Three conclusions drawn from the inverted form were withdrawn: that
\(\Xi\)-boundedness is weaker than \(c_K\)-boundedness; that
\texttt{conj:K-refined} could be weakened to an \(O(1)\) statement without
loss; and that the measured \(c_K\approx1\) was direct evidence for
\(\Xi=O(1)\).

\subsection{What the correction does to the headline}

It reverses it. The corrected reading is delivered by
\texttt{prop:s44-linfty-route} (line~9242), which was re-derived by hand
before merging: routed through \(\norm{w}_{L^\infty}\) --- the route needing no
measure information at all --- \texttt{target:band-absorption} requires
\[
  c_K^{\mathrm{band}} \;\le\; \frac{C_1}{\Lfac},
\]
the \emph{exact analogue} of \texttt{conj:K-refined} on the band. So band
absorption is \textbf{not weaker}; the logarithm is \textbf{not an artefact}
but the price of that route; and in this narrow sense the document's ``same
analytic type'' language is vindicated --- precisely the opposite of what the
first draft reported. The \(O(1)\) alternative exists only in the
volume-normalised statistic, and that statistic is strictly stronger.

The mechanism of the proof is worth one sentence, because it shows where the
logarithm enters and leaves: the region measure \(\abs{R}\) cancels exactly
against \(\iint_R w\psi^2\le\norm{w}_{L^\infty}\abs{R}\); the enstrophy budget
contributes \(\Lfac M^{1/2}\); \(M^{-3/2}=(\rstar)^3\); and the hypothesised
\(C_1/\Lfac\) is exactly what cancels the logarithm the budget carries,
landing the term in the class \texttt{prop:second-rung} absorbs.

\begin{figure}[ht]
\centering
\resizebox{\textwidth}{!}{%
\begin{tikzpicture}[
  >={Latex[length=2mm]},
  bx/.style={draw, rounded corners=2pt, font=\scriptsize, align=center,
             inner sep=3.5pt, text width=34mm, minimum height=11mm},
  op/.style={bx, draw=openblue, thick, fill=panelgrey},
  ar/.style={->, thick},
  no/.style={->, thick, impossiblered}
]
\node[op] (ck) at (0,0)
  {\texttt{conj:K-refined}\\ \(c_K\le C/\Lfac\) on \(E\)\\ \OPEN};
\node[op] (band) at (7.4,0)
  {\texttt{target:band-absorption}\\ \eqref{eq:bandY}, \eqref{eq:bandD} on
   \(\Band\)\\ \OPEN};

\draw[no] (ck) to[bend left=18]
  node[above, font=\scriptsize, text=impossiblered]
  {\texttt{prop:s44-a-fails}} (band);
\draw[no] (band) to[bend left=18]
  node[below, font=\scriptsize, text=impossiblered]
  {\texttt{prop:s44-b-fails}} (ck);
\node[font=\Large, text=impossiblered] at (3.7,0.62) {\(\nRightarrow\)};
\node[font=\Large, text=impossiblered] at (3.7,-0.62) {\(\nLeftarrow\)};

\node[op] (lc) at (3.7,-3.2)
  {\textbf{(LC)}, \texttt{conj:lc-s44}\\
   the \(c_K\) hypothesis verbatim,\\ on \(\{m\ge M/16\}\)\\ \OPEN};
\draw[ar, provedgreen] (lc.west) to[out=170,in=-90]
  node[left, font=\scriptsize, text=provedgreen, pos=0.6]
  {covers the \(Y\)-half} (ck.south);
\draw[ar, provedgreen] (lc.east) to[out=10,in=-90]
  node[right, font=\scriptsize, text=provedgreen, pos=0.6]
  {covers the \(Y\)-half} (band.south);

\node[bx, draw=impossiblered, text width=44mm] (dhalf) at (11.6,-3.2)
  {The \(D\)-half \eqref{eq:bandD} is covered by \textbf{nothing}:
   \texttt{rem:s44-bandD} shows it contains no product of densities and no
   correlation at all};

\node[font=\scriptsize\itshape, text=rulegrey, align=center, text width=96mm]
  at (4.4,-5.3)
  {The scope of the two crossed arrows: no implication follows from
   (A1)--(A3). This is \emph{not} independence for Navier--Stokes solutions
   (\texttt{rem:s44-scope}, line~8855).};
\end{tikzpicture}}
\caption[The \Sn{44} relation map]{What \Sn{44} settled. Neither conjecture
implies the other over the a priori structure they are written against; a
single new hypothesis covers the \(Y\)-half of both; and the \(D\)-half of the
band target is covered by no hypothesis in the family.}
\label{fig:s44-map}
\end{figure}

\section{The common target, and its exact position}
\label{sec:lc}

A common hypothesis does exist, and the better of the two candidates was found
second.

Round one proposed \texttt{conj:jle-s44} (JLE, line~9512), built on the
volume-normalised excess \(\Xi\). It is strong: (JLE-1) yields \(K\)-absorption
outright, with no angular term, no \(\norm{w}_{L^\infty}\) and no annulus
residual, and it also yields \eqref{eq:bandY}. But it degrades like
\(\abs{Q(\rstar)}/\abs{S}\) under concentration on a small set \(S\), which is
exactly the regime near a singularity.

Round two produced the better one.

\begin{statement}[(LC) --- \texttt{conj:lc-s44}, line~9445]
\label{st:lc}
\OPEN{}. (LC) is \emph{verbatim} \texttt{conj:K-refined} --- same
normalisation, same \(C/\Lfac\) decay --- asserted on \(\{m\ge M/16\}\) instead
of \(\{m\ge M/4\}\). That region contains both \(E\) and \(\Band\), and the
integrand is non-negative, so \texttt{prop:s44-lc-implies} (line~9461) bounds
both \(Y\)-halves \emph{with no volume hypothesis}.
\end{statement}

Its position is easy to misstate and the document is emphatic about it.
\texttt{rem:s44-lc-reading} (line~9486): (LC) does \textbf{not} imply
\texttt{conj:K-refined} --- correlation ratios are not monotone under
restriction --- and the converse fails too. (LC) supplies the
\emph{conclusion} both conjectures were designed to give, not either
conjecture as a statement. It is ``a genuinely new hypothesis of an old
form''. It is weaker than JLE in the regime that matters; it is \emph{not}
weaker than \texttt{conj:K-refined}.

\begin{obsv}[The second overstatement in the summaries]
\label{obs:lc-not-weaker}
\texttt{HANDOVER-S46-OPUS.md} described (LC) as ``a common \emph{weaker}
target'' than \texttt{conj:K-refined}. By \texttt{rem:s44-lc-reading} it is
neither weaker nor stronger; the two do not compare.
\texttt{rem:s47-handover-corrigenda} records the correction. This is the
second of the three overstatements \Sn{47} caught in work one session old, and
like the first --- the audit's missing eighth defect --- it runs in the
programme's own favour.
\end{obsv}

\section{The half nothing covers}
\label{sec:dhalf}

\texttt{rem:s44-bandD} is a correction to the surrounding prose and is easy to
overlook. Equation~\eqref{eq:bandD}, the \(D\)-half of the closure target,
contains \emph{no product of two densities and no correlation at all}: it is a
bare level non-concentration statement for \(\nu\abs{\nabla^2\ehat}^2\psi^2\).
No decorrelation hypothesis in this family bears on it.

So ``same analytic type as \texttt{conj:K-refined}'' is accurate for the
\(Y\)-half and \emph{inaccurate} for the \(D\)-half. Future sessions should not
assume the \(D\)-half comes along for free. Nor is the constant condition
\(CC_1\le\delta\nu\) covered by any of the hypotheses
(\texttt{rem:s44-channel-caveat}).

\section{What was not measured, and what \texorpdfstring{\Sn{45}}{S45} then measured}
\label{sec:s44-numerics}

\texttt{rem:s44-numerics} (line~9396) names three quantities that bear on all
of this and had not been measured in the form needed, and it warns that
\Sn{42}--\Sn{43}'s volume fractions were taken over the whole periodic box
rather than over \(Q(\rstar)\), so they are motivation and \emph{not evidence}:

\begin{enumerate}[label=(\roman*),leftmargin=2.4em]
\item the volume fractions \(\abs{E}/\abs{Q(\rstar)}\) and
  \(\abs{\Band}/\abs{Q(\rstar)}\) \emph{at scale} \(\rstar\) --- the
  hypotheses (NB-E) and (NB);
\item the band correlation constant \(c_K^{\mathrm{band}}\);
\item the \(\Lfac\)-dependence of both correlation constants, since
  \texttt{prop:s44-linfty-route} needs \emph{decay}, not boundedness, and
  \Sn{34}--\Sn{35} explicitly could not distinguish \(c_K=O(1)\) from
  \(c_K\le C/\Lfac\).
\end{enumerate}

\Sn{45} measured all three. The results are in Chapter~\ref{ch:layer4}: the
\(\Lfac\)-slope of \(\log c_K\) came back at \(-0.196\) (\(r^2=0.05\)) on one
pool and \(+0.088\) (\(r^2=0.02\)) across all twelve runs, against the \(-1\)
a clean \(C/\Lfac\) law requires, over \(\Lfac\in[1.35,4.68]\); and (NB) failed
outright in one run, where the band was \emph{empty} inside \(B_{\rstar}\).
Chapter~\ref{ch:numerical-lessons} explains why neither number is evidence
about either conjecture, and why that was knowable in advance.

\section{Net effect}
\label{sec:s44-net}

\begin{record}
The open-item count did \textbf{not} fall. Both conjectures remain \OPEN{},
and no implication between them can be derived from (A1)--(A3). What improved
is the targeting: one sharply-stated hypothesis (LC) covers both \(Y\)-halves
instead of two separate ones; the \(D\)-half is explicitly flagged as covered
by nothing; and the \(O(1)\) versus \(C/\Lfac\) confusion is resolved --- it
was a route distinction, not a strength distinction.
\end{record}

\begin{obsv}[The value of checking a resemblance]
\label{obs:resemblance}
Three separate places in the document had described the two conjectures as
being of the same analytic type, and one had suggested that proving either
would deliver the other. That belief was in the record for three sessions and
would have justified spending effort on whichever looked easier.
\Sn{44} cost one session and established that the effort would not have
transferred.

\smallskip
This is the third instance in this book of the same shape: a session
commissioned in the hope of a reduction returns a proof that the reduction
does not exist, and the programme is better off for knowing. \Sn{41} is the
first, \Sn{38} the second in its own way, and \Sn{44} the third. In none of
the three did the mathematics get easier; in all three the map got more
accurate, and in two of the three the correction was in the direction that
cost the programme something.
\end{obsv}
